\chapter{The robot and the system}
\label{ch:u-robot}

\section{What it does}

The X30 walks routes taught by an operator and inspects the equipment at their stop points:

\begin{itemize}
  \item \textbf{Teach once, repeat}: an operator walks the robot along the route with the H12 and
    marks where it stops. The robot repeats the route as often as asked, on its own or on a
    schedule. There is no map: the route is the robot's memory of the walk.
  \item \textbf{Inspect at every stop}: the thermal camera looks at each piece of equipment taught
    there and measures its temperatures; the acoustic camera listens for partial discharge.
    Every finding has a severity, from normal to critical.
  \item \textbf{Record}: during every patrol the cameras' video, their raw thermal frames and the
    acoustic data are recorded on the robot.
  \item \textbf{Report}: NETRA, the command center, shows the robot live, keeps every patrol,
    every inspection with its images, and the trend of every piece of equipment.
\end{itemize}

\section{Its parts}

\begin{center}
\begin{tabularx}{\linewidth}{@{}>{\bfseries}p{0.34\linewidth}X@{}}
  \toprule
  Part & What it is for \\
  \midrule
  Jueying X30 Pro & the four-legged robot, Deep Robotics \\
  HIKMICRO PTZ camera & thermal and visible camera that pans, tilts and zooms \\
  Hertzino HZ-FA-371 & acoustic camera for partial discharge \\
  Four LiDARs & see around the body: obstacles, the robot's own position \\
  Skydroid H12 & the handheld with the X30Controller app \\
  NETRA & the command center on a laptop or server \\
  \bottomrule
\end{tabularx}
\end{center}

\section{Its start spot}

The robot knows where it is relative to the spot it was started on, usually its dock. Start it
there every time, facing the same way: every route was taught from that spot, and the robot
repeats it from there. Markers on the walls (tags), where the site has them, correct the small
drift that grows with every metre walked.
